% =====================================================================
\section{Related Work}
\label{sec:related}

\noindent\textbf{Synthetic data for correspondence.}
FlyingChairs/FlyingThings \cite{mayer2016FlyingThingsDataset}, Sintel
\cite{sintel}, Kubric/MOVi \cite{greff2021kubric}, and PointOdyssey
\cite{pointodyssey} established synthetic pre-training as the default for
dense matching; AutoFlow \cite{sun2021autoflow} and DFlow \cite{im2022dflow}
\emph{learn} data parameters with training in the loop. Closest to our
premise, Mayer~\etal~\cite{mayer2018what} showed that matching target
displacement statistics improves flow training. That result is undirected;
our finding is the part it could not see. The match is
\emph{directional}: coverage of the target's motion is the half that predicts
transfer, and a pretrained backbone is what lets a model use it. A large, active
line of work improves the \emph{rendering} side of these pipelines (photorealism,
lighting, scale, speed) \cite{greff2021kubric,pointodyssey}; this is complementary:
we take the renderer as given and instead supply a target-conditioned signal for
\emph{what} to generate, which is why the same rule guides two structurally
different generators (Sect.~\ref{sec:interv}).
Procedural abstraction \cite{Baradad2021LookingAtNoise,white2009mandelbulb}
supports our use of non-photorealistic, motion-controllable generators.

\noindent\textbf{Transferability estimation.}
LEEP \cite{nguyen2020leep}, LogME \cite{you2021logme}, NCE
\cite{tran2019transferability}, H-score \cite{bao2019information}, and OTCE
\cite{tan2021otce} predict transfer from model outputs or target labels,
almost exclusively for classification. OTDD \cite{alvarez2020geometric} is,
like ours, a fit-free dataset distance, but symmetric and label-dependent---its
label-to-label optimal transport assumes a discrete class space, which dense
correspondence lacks, so it does not directly apply here. Our sliced-$W_2$ baseline
is itself a (label-free) optimal-transport distance, so the OT family \emph{is}
represented among our symmetric comparisons. None of these methods address dense
correspondence, and none are directional; our experiments show the symmetric,
direction-averaging choice---not the discrepancy itself---is why such signals wash
out.

\noindent\textbf{Precision/recall divergences.}
Directional two-sample diagnostics are established in generative-model
evaluation \cite{sajjadi2018assessing_precision_recall,
kynkaanniemi2019improved_pr,naeem2020reliable,alaa2022faithful}. We adopt the
same decomposition in motion space and make it \emph{predictive and
directional} for transfer, which to our knowledge is new.

\noindent\textbf{Context-conditional data value.}
Sorscher~\etal~\cite{sorscher2022beyond} showed the value of easy vs.\ hard
examples flips with data abundance; pretraining-data selection methods such
as DSIR \cite{xie2023data} and DoReMi \cite{xie2023doremi} match
distributions symmetrically. Our result is a new instance of
context-conditional data value: a source's worth is set by how well its motion
\emph{covers} a given target's---a directed quantity, not the symmetric
distribution match those methods optimize. On the theory side, asymmetric transfer exponents
\cite{hanneke2019value} and covariate-shift support conditions anticipate the
headline---coverage of the target's support predicts transfer---but not how cheaply
it can be measured (a crude motion histogram, no training) or that off-target mass
is nearly free. Our dataset-size controls (Sect.~\ref{sec:law}) rule out a trivial
``more data'' explanation.

\noindent\textbf{Architectures and backbones.}
We span semantic matching (CATs++ \cite{cho2022cats++}), recurrent flow (RAFT
\cite{teed2020raft}), and pyramid dense matching (GLU-Net
\cite{truong2020glunet}); appearance descriptors use DINOv3
\cite{simeoni2025dinov3}, whose strong zero-shot matching
\cite{sd_dino_tale_of_two_features} motivates appearance as a serious
competing hypothesis---which our controls reject for within-context dataset
ranking.

